\begin{abstract}
Alzheimer's disease (AD) involves many biological processes at once, so treatment is expected to need combinations of drugs, often ones already approved for other conditions. Computationally, choosing $k$ drugs from $n$ candidates is a search over $\binom{n}{k}$ combinations. We follow a network-medicine rule for good combinations: each drug's protein targets should lie close to the AD disease proteins, while the drugs' target sets should stay separated from one another. We encode this rule as a single integer-valued score (``energy'') over every combination, in a quadratic form suited to both annealing and QAOA, and search for its minimum in three ways: (1)~An exhaustive GPU search checks every combination and computes the true minimum, which serves as ground truth. (2)~A simulated-annealing search is implemented as sequential C++, as an AVX2SIMD kernel written in assembly, and as a CUDA kernel. All three implementations use the same per-chain random streams and fixed-point integer arithmetic, ideally returning identical solutions. Any timing difference therefore reflects the underlying hardware rather than the algorithm. We report throughput and speedup over the sequential baseline as $n$, $k$, and the number of parallel chains grow. (3)~ Moving away from traditional computing, we further implement a Quantum Approximate Optimization Algorithm (QAOA) running on a simulated quantum computer implemented via NVIDIA cuQuantum. In benchmarking, we measure how close it gets to the exact minimum, along with its simulation time and memory, which grow exponentially with $n$. As a limitation, we show that the results are for computational hypotheses for later experimental testing, not treatment recommendations, and because the quantum computer is simulated on classical hardware, we strictly cannot conclude quantum speedup.\end{abstract}